% ============================================================================
% Fichier  : partie6/P06_C31_mesurer_incertitude_strategie.tex
% Titre    : Mesurer l'Incertitude et la Stratégie
% Auteur   : MINKA MI NGUIDJOÏ Thierry Emmanuel
% Version  : 1.0 -- novembre 2026
% Base   : ACSI C22, repris et développé pour INF3226
% ============================================================================
% Sources : A1 (hal-05701270) : espace Omega_q, distribution p_q, conditions
% C1 à C4 de la réponse NEK, métriques H_q, U_q, mu_q, rho_q, D_KL, quadrant,
% cas limites, SUI (modèle logistique par répondant, score d'enjeu exogène,
% trois réserves), protocole pré-enregistré à trois bras (G1, G2, G3).
% Décision du 06/11/2026 (point 12 de l'auteur) : les travaux non publiés
% (A2 à A5) sont retirés ; les prolongements s'appuient sur des travaux
% publiés et vérifiables (Fellingham et Newman 1985, Spence 1973, JCGM 100).
% ATTENTION : la Proposition 1 (décomposition de l'entropie) de la
% prépublication A1, version HAL comprise, porte un terme supplémentaire
% « + U_q » ; il conduit, pour la distribution uniforme, à H > log2 3, ce qui
% est impossible. Le cours donne l'identité exacte (propriété de groupement
% de l'entropie), vérifiée sur les valeurs du cours. Erratum à décider par
% l'auteur pour A1.
% ============================================================================

\chapter{Mesurer l'Incertitude et la Stratégie}
\label{chap:incertitude}
\chaptermark{Incertitude et stratégie}

\epigraph{%
    « Distinguer ce que l'on sait,\\
    ce que l'on a déposé, et ce qui reste à prouver. »%
}{Règle du chapitre}

\niveauChapitre{Recherche}{%
    Maîtrise du Chapitre~\ref{chap:auditologie} (axiomes et piliers) et du
    Chapitre~\ref{chap:conclure} (mesures)%
}

\begin{boxobjectifs}
\begin{itemize}
    \item Approfondir l'instrument trichotomique et ses métriques (A1)~:
          conditions de la réponse \NEK{}, décomposition de l'entropie, cas
          limites.
    \item Comprendre l'indice d'usage stratégique (SUI), ses conditions
          d'identification et ses limites.
    \item Relier l'instrument à des travaux publiés~: théorie des jeux de
          l'audit, signalement, expression de l'incertitude de mesure.
    \item Distinguer ce qui est établi, ce qui est déposé en prépublication
          et ce qui reste à valider.
\end{itemize}
\end{boxobjectifs}

\bigskip

% ============================================================================
\section*{Question directrice : que vaut ce que nous mesurons~?}
% ============================================================================

Le cours a calculé des dizaines de valeurs de $\mu$, $U$ et $H$. Elles
reposent sur un instrument formalisé dans l'article A1~\cite{minka2026nek}.
Ce chapitre en reprend les fondements, puis le relie à des travaux publiés
qui en éclairent les prolongements~: la stratégie des répondants et de
l'auditeur, et l'expression d'une incertitude de mesure. Il se termine par
une règle de lecture~: tout n'a pas le même statut.

\medskip
\textbf{Plan du chapitre.}
\begin{enumerate}
    \item L'instrument trichotomique (section~\ref{sec:c31-instrument}).
    \item La stratégie~: l'indice d'usage stratégique (section~\ref{sec:c31-sui}).
    \item Des travaux publiés pour prolonger l'instrument (section~\ref{sec:c31-travaux}).
    \item Établi, déposé, à valider (section~\ref{sec:c31-statut}).
\end{enumerate}

\separateur

% ============================================================================
\section{L'Instrument Trichotomique}
\label{sec:c31-instrument}
% ============================================================================

\subsection{L'espace des réponses}

Pour chaque question $q$ d'un ensemble $I$, A1 remplace l'espace binaire
par $\Omega_q = \{\OUI{}, \NON{}, \NEK{}\}$. Pour une population de
répondants $R$, la distribution observée est
\[
    p_q(m) = \frac{1}{|R|} \sum_{r \in R} \mathbf{1}[x_{r,q} = m],
    \qquad m \in \Omega_q,
\]
et l'on note $\mu_q = p_q(\OUI{})$ et $U_q = p_q(\NEK{})$. Le formalisme ne
fait aucune hypothèse sur la manière dont les croyances se forment~: les
engagements portent sur l'usage de \NEK{}.

\subsection{Quatre conditions pour \NEK{}}

A1 précise quatre conditions qui, ensemble, distinguent un \NEK{} légitime
d'un mauvais usage. Le cours les a appliquées sans les nommer.

\begin{tableaucours}{L{3.6cm}L{5.0cm}L{3.6cm}}{%
    \tableauHeader{Condition (A1)} &
    \tableauHeader{Sens} &
    \tableauHeader{Dans le cours}
}
C1. Spécificité de l'objet & \NEK{} porte sur l'objet précis de la
question, non sur le thème & Une question par contrôle, atomique
(Chapitre~\ref{chap:decomposition}) \\
C2. Asymétrie d'accès & L'information existe dans l'organisation, mais
le répondant n'y a pas accès & Le document non remis dont on ignore
l'existence (Chapitre~\ref{chap:planifier}) \\
C3. Non-report & \NEK{} n'est pas « je vérifie et je reviens »~: si l'on
peut vérifier dans le délai de la mesure, on répond & Délai écrit pour
produire les pièces (Chapitre~\ref{chap:droit}) \\
C4. Neutralité & \NEK{} n'approuve ni ne désapprouve~: aucun barème ne
doit le rendre équivalent à \OUI{} ou à \NON{} & Pas de « demi-point »
pour \NEK{} \\
\end{tableaucours}
\vspace{-0.8em}
\begin{center}{\small\itshape Tableau~31.1. Les conditions C1 à C4 de la
réponse \NEK{} (d'après A1)}\end{center}

\subsection{Les métriques}

\begin{tableaucours}{cL{4.2cm}L{6.4cm}}{%
    \tableauHeader{Mesure} &
    \tableauHeader{Définition} &
    \tableauHeader{Interprétation}
}
$H_q$ & $-\sum_m p_q(m) \log_2 p_q(m)$ & Dispersion des états de
connaissance des répondants \\
$U_q$ & $p_q(\NEK{})$ & Lacune collective de connaissance \\
$\mu_q$ & $p_q(\OUI{})$ & Taux de conformité~; rejoint le taux binaire
quand $U_q$ est faible \\
$\rho_q$ & $\hat\mu_q^{B} - \mu_q^{T}$ & Biais du choix forcé, du binaire
au trichotomique \\
$D_{\mathrm{KL}}$ & $\sum_m p^{(t+1)}(m) \log_2 \frac{p^{(t+1)}(m)}{p^{(t)}(m)}$
& Coût informationnel de prédire le cycle $t+1$ à partir du cycle $t$ \\
\end{tableaucours}
\vspace{-0.8em}
\begin{center}{\small\itshape Tableau~31.2. Les métriques de l'instrument
(d'après A1)}\end{center}

Deux remarques d'A1 sur la divergence méritent d'être retenues. Elle est
asymétrique~: on la calcule dans le sens du temps, du cycle passé vers le
cycle présent. Et elle exige qu'une modalité observée au cycle $t+1$ l'ait
été aussi au cycle $t$~; sinon, on lisse les fréquences des petits
échantillons.

\subsection{Décomposer l'entropie}

L'entropie trichotomique mêle deux sources de dispersion~: l'ignorance des
répondants, et le désaccord entre ceux qui savent. La propriété de
groupement de l'entropie~\cite{shannon1948} les sépare exactement. En
notant $H_b(x) = -x\log_2 x - (1-x)\log_2(1-x)$ l'entropie binaire~:
\[
    H_q \;=\; H_b(U_q) \;+\; (1 - U_q)\, H_b\!\left(\frac{\mu_q}{1 - U_q}\right).
\]
Le premier terme mesure la part de l'ignorance~; le second, pondéré par la
part des répondants informés, mesure leur désaccord. Une grille binaire
n'observe que le second, et seulement après avoir forcé les ignorants à
choisir.

\begin{boxlecture}{Deux entropies, deux causes}
\begin{tableaucours}{L{2.4cm}cccccc}{%
    \tableauHeader{Domaine} &
    \tableauHeader{O/N/K} & \tableauHeader{$U$} &
    \tableauHeader{$H_b(U)$} & \tableauHeader{Désaccord} &
    \tableauHeader{$H$} & \tableauHeader{Part ignorance}
}
D, mai~2026 & 13/2/5 & 0,25 & 0,81 & 0,42 & 1,24 & 66~\% \\
E, mai~2026 & 8/6/6 & 0,30 & 0,88 & 0,69 & 1,57 & 56~\% \\
B, mai~2026 & 8/3/1 & 0,08 & 0,41 & 0,77 & 1,19 & 35~\% \\
\end{tableaucours}

\medskip
\textbf{Constatation.} Les trois domaines ont une entropie supérieure à
1~bit, pour des raisons différentes. Dans D, les deux tiers de
l'entropie viennent de l'ignorance~; dans B, près des deux tiers viennent
du désaccord entre répondants informés.

\medskip
\textbf{Analyse} (confiance élevée sur le calcul). La décomposition dit
quoi faire~: dans D, chercher la preuve~; dans B, chercher pourquoi des
répondants qui savent ne disent pas la même chose (un site, une
procédure, une période). C'est la lecture « $H$ élevé, $U$ faible » du
Chapitre~\ref{chap:conclure}, rendue quantitative.

\medskip
\hfill\lecture{0{,}65}{0{,}25}{1{,}24}
\end{boxlecture}

\subsection{Cas limites}

A1 énonce quatre cas limites. L'accord unanime sur \OUI{} donne $H = 0$,
$U = 0$, $\mu = 1$~; sur \NON{}, $H = 0$, $U = 0$, $\mu = 0$. L'ignorance
unanime donne $H = 0$ et $U = 1$~: entropie nulle, mais révélation
maximale, puisque chacun sait qu'il ne sait pas. La distribution uniforme
donne $H = \log_2 3$ et $U = \mu = 1/3$. D'où la règle du cours~: on ne
lit jamais $H$ seul.

% ============================================================================
\section{La Stratégie : l'Indice d'Usage Stratégique}
\label{sec:c31-sui}
% ============================================================================

\subsection{Le problème d'identification}

Un \NEK{} est \emph{authentique} s'il traduit une vraie limite de
connaissance~; il est \emph{stratégique} si le répondant sait, mais évite
les conséquences d'une réponse nette. Pour une réponse isolée, les deux
sont indiscernables. A1 propose de les distinguer au niveau du
\emph{profil}~: l'ignorance authentique suit l'accès à l'information, qui
ne dépend pas, en moyenne, de l'enjeu de la question~; l'ignorance
stratégique se concentre sur les questions à fort enjeu.

\subsection{La construction}

Pour chaque question $q$, un score d'enjeu $s_q \in [0, 1]$ est fixé
\emph{avant} la collecte par un panel d'experts indépendants. Pour chaque
répondant $r$, A1 estime un modèle logistique
\[
    \log \frac{\Pr[\NEK{}_{r,q} = 1]}{1 - \Pr[\NEK{}_{r,q} = 1]}
    \;=\; \alpha_r + \beta_r\, s_q + \gamma_r^{\top} \mathbf{z}_q,
\]
où $\mathbf{z}_q$ contient des contrôles (domaine, longueur, difficulté
technique), et définit $\mathrm{SUI}_r = \hat\beta_r$. Un indice proche de
zéro est compatible avec un usage authentique~; un indice positif et
significatif signale un usage stratégique. L'ordonnée à l'origine
$\alpha_r$ absorbe la propension de chacun à répondre \NEK{}~: un
répondant qui en use beaucoup, mais partout, n'est pas stratégique.

\begin{boxscenario}{Opération LIZBIZ. Deux contextes, mêmes questions}
Au Chapitre~\ref{chap:fraude}, les mêmes dix questions sur le stock de
Douala, posées aux mêmes deux responsables, donnaient 17~\OUI{}, 2~\NON{},
1~\NEK{} en auto-évaluation, et 9~\OUI{}, 3~\NON{}, 8~\NEK{} en entretien
formel. L'écart entre les deux contextes est le signal qu'A0 décrit aussi
pour le domaine~D. Il ne suffit pas à calculer un SUI~: il faudrait des
scores d'enjeu fixés à l'avance par des experts indépendants, et assez de
questions par répondant pour estimer le modèle.
\end{boxscenario}

\begin{boxalert}
\textbf{Trois réserves d'A1, à ne jamais oublier.}
\begin{enumerate}
    \item Un SUI positif signale une corrélation entre \NEK{} et enjeu, qui
          peut avoir d'autres causes~: questions plus difficiles, accès plus
          restreint.
    \item Le score d'enjeu doit être exogène~: s'il est fixé par les
          répondants eux-mêmes, l'indice n'est pas identifié.
    \item Le SUI est un indicateur de population, \emph{pas} un diagnostic
          de faute individuelle~: il n'établit pas l'intention et ne fonde
          aucune sanction.
\end{enumerate}
\end{boxalert}

% ============================================================================
\section{Des Travaux Publiés pour Prolonger l'Instrument}
\label{sec:c31-travaux}
% ============================================================================

\subsection{Le jeu entre l'auditeur et l'audité}

L'idée que l'audit est un jeu, où chacun anticipe la stratégie de l'autre,
est ancienne dans la recherche comptable. J.~Fellingham et D.~Newman ont
modélisé dès 1985 l'interaction entre l'effort de l'auditeur et le
comportement de l'audité~\cite{fellingham1985strategic}~; la théorie des
jeux à information incomplète, formalisée par J.~Harsanyi, en donne le
cadre~\cite{harsanyi1967}. L'axiome~4 d'A0 (interaction stratégique)
s'inscrit dans cette tradition~; le SUI d'A1 en est une mesure empirique.

\subsection{Le signal coûteux}

M.~Spence a montré qu'un signal n'informe que s'il coûte davantage à celui
qui voudrait tromper~\cite{spence1973signaling}. Appliqué à l'audit~: une
réponse \NEK{} ne coûte rien à donner~; elle ne peut donc pas, seule, être
un signal fiable d'ignorance sincère. C'est pourquoi le SUI compare les
réponses selon l'enjeu, plutôt que de les prendre à la lettre.

\subsection{Exprimer une incertitude de mesure}

La métrologie dispose depuis longtemps d'un cadre pour exprimer
l'incertitude d'une mesure~: le \emph{Guide pour l'expression de
l'incertitude de mesure}~\cite{jcgm2008gum}. Il distingue les évaluations
statistiques (de type~A) et les autres (de type~B, fondées sur le jugement
et l'information disponible). Les mesures du cours se rangent dans les
deux~: $\mu$ et $U$ sont des proportions observées, dont on peut donner un
intervalle (type~A)~; les seuils du quadrant sont des conventions déclarées
(type~B). Un rapport rigoureux le dit.

\begin{boxinfo}
\textbf{Un intervalle pour $U$.} Avec 5~\NEK{} sur 20~réponses, $U = 0{,}25$.
L'intervalle de confiance à 95~\% d'une proportion observée sur 20 cas est
large~: approximativement de 0,11 à 0,47 (méthode de Wilson). Classer un
domaine « opaque » sur un seuil de 0,20 avec si peu de réponses est donc
une décision, à motiver, plus qu'une constatation. C'est une raison de
plus pour multiplier les répondants sur les domaines importants.
\end{boxinfo}

% ============================================================================
\section{Établi, Déposé, à Valider}
\label{sec:c31-statut}
% ============================================================================

\begin{tableaucours}{L{3.6cm}L{3.4cm}L{5.2cm}}{%
    \tableauHeader{Élément} &
    \tableauHeader{Statut} &
    \tableauHeader{Ce qui reste à faire}
}
Entropie, divergence, groupement & Établi (théorie de l'information) &
Rien~: résultats classiques \\
Programme de l'auditologie (A0) & Prépublication HAL, 2026 & Discussion
par la communauté~; tests des prédictions \\
Instrument \OUI{}/\NON{}/\NEK{}, métriques, SUI (A1) & Prépublication HAL,
2026~; protocole pré-enregistré & Expérience à trois bras~: groupes
fondation, validation, généralisation, chacun en binaire et en
trichotomique \\
Jeu de l'audit, signal coûteux, incertitude de mesure & Établi (travaux
publiés) & Les appliquer explicitement à l'instrument \NEK{} \\
Cas du cours (LiZbiZ, SANAGA, ANPP) & Fictifs, pédagogiques & Ils
illustrent~; ils ne prouvent rien \\
\end{tableaucours}
\vspace{-0.8em}
\begin{center}{\small\itshape Tableau~31.3. Le statut des éléments
présentés dans le cours}\end{center}

\begin{boiteRemarque}{Pourquoi insister sur le statut~?}
Parce que l'auditologie s'applique à elle-même ce qu'elle demande à
l'auditeur~: distinguer ce qui est établi (\OUI{}), ce qui est réfuté
(\NON{}) et ce qui reste inconnu (\NEK{}). Un étudiant qui cite le SUI
comme un résultat validé commettrait l'erreur que le cours apprend à
éviter~: coter \OUI{} sans preuve.
\end{boiteRemarque}

\begin{boxmetier}{Auditeur SI}
\textbf{Contribuer à la recherche sans trahir la mission}
\begin{enumerate}
    \item Conserver les cotations question par question, avec le contexte
          (entretien, auto-évaluation) et la date.
    \item Garder les mêmes questions et les mêmes critères d'un cycle à
          l'autre, pour que la divergence ait un sens.
    \item Anonymiser avant tout partage, avec l'accord écrit du
          commanditaire.
    \item Ne jamais utiliser un indicateur de recherche, comme le SUI, pour
          mettre en cause une personne.
\end{enumerate}
\end{boxmetier}

% ============================================================================

\separateur

\begin{boxresume}
\textbf{Ce qu'il faut retenir de ce chapitre~:}
\begin{itemize}
    \item \textbf{\NEK{} obéit à quatre conditions} (A1)~: spécificité de
          l'objet, asymétrie d'accès, non-report, neutralité.
    \item \textbf{L'entropie se décompose} en une part d'ignorance,
          $H_b(U)$, et une part de désaccord entre répondants informés~;
          elle dit s'il faut chercher la preuve ou la cause du désaccord.
    \item \textbf{Le SUI} repère un usage stratégique de \NEK{} au niveau
          d'une population, à partir d'enjeux fixés à l'avance~; il ne
          désigne aucun coupable.
    \item \textbf{Des travaux publiés éclairent l'instrument}~: le jeu de
          l'audit, le signal coûteux, l'expression de l'incertitude de
          mesure (intervalle autour de $U$).
    \item \textbf{Établi, déposé, à valider}~: le cours applique à ses
          propres sources la distinction \OUI{}, \NON{}, \NEK{}.
\end{itemize}
\end{boxresume}

\begin{boxinfo}
\textbf{Points de contrôle de l'élève.} Je sais~: énoncer les quatre
conditions de \NEK{}~; calculer $\mu$, $U$, $H$ et décomposer $H$~; dire ce
que le SUI mesure et ne mesure pas~; donner un intervalle autour de $U$~;
distinguer incertitude de type~A et de type~B~; donner le statut de
chaque source.
\end{boxinfo}

\bigskip

\begin{boxexo}[title={\ding{45}\quad Exercice~31.1. Quelle condition~? \niveaubase}]
Indiquez la condition C1 à C4 que chaque situation méconnaît.
\begin{enumerate}[(a)]
    \item Un chef d'agence répond \NEK{} à toute la section « sécurité »,
          qu'il ne connaît pas.
    \item Un répondant coche \NEK{} en attendant de consulter un dossier
          qu'il a sur son bureau.
    \item Une grille compte \NEK{} pour un demi-point de conformité.
    \item Un DSI répond \NEK{} sur un contrôle qu'il administre lui-même.
\end{enumerate}
\end{boxexo}

\begin{boxexo}[title={\ding{45}\quad Exercice~31.2. Décomposer l'entropie \niveaumoyen}]
\begin{enumerate}[(a)]
    \item Décomposez l'entropie du domaine~E de LiZbiZ en novembre~2026
          (8/4/8) et comparez la part d'ignorance à celle de mai.
    \item Vérifiez l'identité de groupement sur la distribution uniforme
          (1/1/1).
    \item Décomposez le cycle achats-paiements après tests (14/6/0) et
          interprétez.
\end{enumerate}
\end{boxexo}

\begin{boxexo}[title={\ding{45}\quad Exercice~31.3. Préparer un sujet de mémoire \niveauexpert}]
Choisissez une prédiction de la matrice de falsifiabilité
(Chapitre~\ref{chap:auditologie}) ou une hypothèse d'A1.
\begin{enumerate}[(a)]
    \item Formulez la question de recherche et l'hypothèse.
    \item Décrivez le terrain, les données, et l'accord à obtenir.
    \item Précisez ce qui réfuterait l'hypothèse, et ce qui serait révisé.
    \item Indiquez le statut de chaque source citée~: publiée et évaluée, ou
          déposée en prépublication.
\end{enumerate}
\end{boxexo}

\bigskip

\begin{boxexo}[title={\ding{45}\quad Exercice~31.4. Un intervalle pour U \niveaumoyen}]
En novembre~2026, le domaine~E de LiZbiZ compte 8~\NEK{} sur 20~réponses.
Calculez l'intervalle de Wilson à 95~\% de $U$ et dites si l'on peut
affirmer que $U$ a baissé par rapport à un cycle où il valait 0,60 sur
20~réponses.
\end{boxexo}

\begin{boxexo}[title={\ding{45}\quad Exercice~31.5. Type A ou type B~? \niveaubase}]
Classez selon le GUM~: (a) l'intervalle autour de $U$~; (b) le seuil de
0,20 qui classe un domaine « opaque »~; (c) le seuil de 0,8 du ratio
d'impact~; (d) l'écart-type d'un taux d'erreur estimé par sondage~; (e)
le jugement d'un expert sur la fiabilité d'un journal.
\end{boxexo}

\begin{boxexo}[title={\ding{45}\quad Exercice~31.6. Rendre le signal coûteux \niveaumoyen}]
Une réponse \NEK{} ne coûte rien. Proposez deux règles de l'instrument qui
rendent plus coûteux un \NEK{} de dissimulation qu'un \NEK{} sincère, en
vous appuyant sur Spence et sur les quatre conditions.
\end{boxexo}

\begin{boxexo}[title={\ding{45}\quad Exercice~31.7. Un jeu d'inspection \niveaumoyen}]
Un agent gagne 10 en commettant une irrégularité non détectée et perd 30 si
elle est détectée lors d'un contrôle. Le contrôle détecte toujours
l'irrégularité. Avec quelle probabilité de contrôle l'agent devient-il
indifférent~? Que se passe-t-il au-dessus et au-dessous~? Quelle leçon pour
la planification de l'audit~?
\end{boxexo}

\begin{boxexo}[title={\ding{45}\quad Exercice~31.8. Le statut des sources \niveaubase}]
Donnez le statut de chaque source (établie, déposée, à valider, non
citable)~: l'ISO~31000~; l'article A1~; Spence (1973)~; un billet de
blogue non signé~; un mémoire de master non soutenu~; le GUM.
\end{boxexo}

\begin{boxexo}[title={\ding{45}\quad Exercice~31.9. Écrire l'incertitude \niveauexpert}]
Rédigez le paragraphe du rapport qui présente le domaine~E de LiZbiZ en
novembre~2026 (8/4/8), avec ses mesures, l'intervalle de $U$, la part
d'ignorance de $H$ et le caractère conventionnel des seuils, sans
technicité inutile pour le comité d'audit.
\end{boxexo}

\begin{boxexo}[title={\ding{45}\quad Exercice~31.10. Concevoir une expérience \niveauexpert}]
Concevez une expérience qui compare une grille binaire et une grille
trichotomique sur le même terrain~: groupes, affectation, mesures,
hypothèse, condition de réfutation, et précautions éthiques.
\end{boxexo}

\bigskip

\begin{boxinfo}
\textbf{Pour aller plus loin.}
\begin{itemize}
    \item L'article A1~\cite{minka2026nek}.
    \item T.~Cover et J.~Thomas, \emph{Elements of Information
          Theory}~\cite{cover2006elements}.
    \item Le GUM~\cite{jcgm2008gum}~; M.~Spence~\cite{spence1973signaling}~;
          J.~Fellingham et D.~Newman~\cite{fellingham1985strategic}.
\end{itemize}
\end{boxinfo}

\begin{boxinfo}
\textbf{Conclusion du cours.} Auditer, c'est réduire l'incertitude d'un
tiers sur un système qu'il ne peut pas observer seul. Le métier en donne
les gestes~; l'auditologie en cherche les raisons, et accepte d'être
réfutée. Les annexes rassemblent le glossaire, les outils et les modèles
pour la pratique.
\end{boxinfo}
